\chapter{Medical Abortion}

\section*{Overview}
Medical abortion is the termination of a pregnancy using medications rather than surgery. It is effective for early pregnancy, generally up to about 10-11 weeks of gestation, and is performed under clinician supervision.

\section*{Alternative Names}
Abortion pill; medication abortion; nonsurgical abortion; mifepristone-misoprostol abortion; RU-486

\section*{Principle}
A combination of two medications is used: mifepristone blocks the action of progesterone, which is needed to maintain the pregnancy, and misoprostol, taken 24-48 hours later, causes uterine contractions that expel the pregnancy. The process produces bleeding and cramping similar to a heavy period.

\section*{History}
Mifepristone (RU-486) was first approved in France in 1988 and in the United States in 2000. The mifepristone-misoprostol regimen has since become the standard medical method for early abortion worldwide.

\section*{Why This Test or Procedure Is Ordered}
Chosen by a person who wishes to end an early pregnancy and prefers to avoid surgery. It is also used for incomplete or missed miscarriage management. The clinician must confirm gestational age before proceeding.

\section*{Is This a Primary or Secondary Test or Procedure}
A primary method of early pregnancy termination, offering a safe alternative to surgical abortion for eligible patients.

\section*{How This Test or Procedure Is Performed}
Mifepristone is taken by mouth in the clinic or at home, followed 24-48 hours later by misoprostol, which may be taken by mouth or placed in the cheek, under the tongue, or in the vagina. Bleeding and cramping typically begin within hours of the misoprostol dose, and a follow-up visit confirms that the abortion is complete.

\section*{What Preparation Is Needed}
Pregnancy dating by ultrasound or last menstrual period, counseling about options, and a plan for follow-up. Rh-negative patients may need Rh immune globulin, and arrangements for pain control and after-hours contact should be made.

\section*{Contraindications and Precautions}
Medical abortion is not used in pregnancies beyond the approved gestational age, in ectopic pregnancy, or when there is an allergy to the medications. Severe anemia, bleeding disorders, and long-term corticosteroid or anticoagulant use require special caution.

\section*{Risks}
\begin{itemize}
  \item Incomplete abortion requiring surgical completion
  \item Heavy or prolonged bleeding
  \item Infection
  \item Persistent pain, nausea, or diarrhea from the medications
\end{itemize}

\section*{What Happens After the Test or Procedure}
Cramping and bleeding usually resolve within 1-2 weeks, and a follow-up visit or home test confirms the pregnancy has ended. Normal activities may resume once bleeding settles, but intercourse and tampon use should be avoided until the bleeding stops.

\section*{Understanding Your Results}
The procedure is successful when bleeding occurs and a follow-up evaluation shows no remaining pregnancy tissue. Symptoms of pregnancy, such as nausea and breast tenderness, usually resolve within days.

\section*{Normal Results}
Complete expulsion of the pregnancy with return of a normal menstrual period, usually within 4-6 weeks, and no need for further intervention.

\section*{Follow-up Tests and Procedures}
A follow-up visit 1-2 weeks later with ultrasound or a pregnancy test to confirm completion. Surgical evacuation is performed if the abortion is incomplete or bleeding is excessive.

\section*{References}
\begin{enumerate}
  \item MedlinePlus. Medical Abortion. U.S. National Library of Medicine; 2025.
  \item Current Medical Diagnosis and Treatment. McGraw-Hill; 2025.
\end{enumerate}

\clearpage
